Let $A : \mathfrak{R}_1 \to \mathfrak{R}_2$ be a linear transformation. Let $(e_1, e_2, \ldots, e_{n_1})$ and $(f_1, f_2, \ldots, f_{n_2})$ be bases in these spaces, and let $(e_1^*, e_2^*, \ldots, e_{n_1}^*)$ and $(f_1^*, f_2^*, \ldots, f_{n_2}^*)$ be the complementary dual bases in $\mathfrak{R}_1^*$ and $\mathfrak{R}_2^*$ respectively:

$$e_i^*(e_j) = \delta_{ij}, \quad f_p^*(f_q) = \delta_{pq}.$$

Suppose $(\alpha) = (\alpha_{ip})$ is the matrix of $A$ relative to the bases $(e_i)$ of $\mathfrak{R}_1$ and $(f_p)$ of $\mathfrak{R}_2$, and $(\beta) = (\beta_{qj})$ is the matrix of the transpose $A^*$ relative to the complementary bases $(e_j^*)$ of $\mathfrak{R}_1^*$ and $(f_q^*)$ of $\mathfrak{R}_2^*$. Then

$$\beta_{qj} = \alpha_{jq}.$$

Thus the matrix of $A^*$ (considering $\mathfrak{R}_i^*$ as left vector spaces) is the transposed matrix $(\alpha)'$ of the matrix $(\alpha)$ of $A$.
